\subsection{Image Retrieval}
\label{sec:fp-res-image}
Table~\ref{tab:fp-image} summarises the image track and Fig.~\ref{fig:fp-image-heatmap} breaks Recall@1 down by transformation. Three groups separate clearly. The copy-detection descriptors trained on DISC21 lead: at the headline operating point (pair-level FPR $10^{-7}$) the ISC21 winner detects \val{fp:image/ISC21-1st/TPR@pair@1e-07/pct}\% of attacked queries and SSCD-R50 \val{fp:image/SSCD-mixup/TPR@pair@1e-07/pct}\%. The best product-eligible methods are the general-purpose DINOv2 embeddings, at \val{fp:image/DINOv2-B/TPR@pair@1e-07/pct}\% (DINOv2-B) and \val{fp:image/DINOv2-S/TPR@pair@1e-07/pct}\% (DINOv2-S), and the perceptual hashes follow, led by PDQ with dihedral matching at \val{fp:image/PDQ-dihedral/TPR@pair@1e-07/pct}\% and plain PDQ at \val{fp:image/PDQ/TPR@pair@1e-07/pct}\%. The ISCC Image-Code reaches \val{fp:image/ISCC-Image-64/TPR@pair@1e-07/pct}\% at its default 64 bits and \val{fp:image/ISCC-Image-256/TPR@pair@1e-07/pct}\% at 256 bits. All paired differences in Table~\ref{tab:fp-paired} are significant after Holm correction.

\begin{figure*}[!t]
\centering
\includegraphics[width=\textwidth]{fp_fig2_image_heatmap}
\caption{Recall@1 of every image method under every transformation, on the ABO registry of \val{fp:count/image/nref} images. Rows are grouped by family (compression, scaling, cropping, geometry, photometric, overlay, chains). Method names are coloured by class; \dag~marks reference-tier methods.}
\label{fig:fp-image-heatmap}
\end{figure*}

\begin{table*}[!t]
\centering
\footnotesize
\caption{Image retrieval over all attacks (ABO; registry of registered originals and distractors). TPR counts a query only if its top-1 is the true asset and the score passes a threshold fixed on held-out calibration negatives: query-level FPR 1\% at this registry size, or pair-level FPR $10^{-7}$ (95\% interval from a bootstrap over sources, which redraws the calibration negatives and so the threshold). Hard: rate at which another photograph of a registered product binds to that product at the 1\% threshold.}
\label{tab:fp-image}
\setlength{\tabcolsep}{3pt}
\begin{tabular}{lcccccc}
\toprule
Method & R@1 & $\mu$AP & TPR$_{q,1\%}$ & TPR$_{p,10^{-7}}$ & Hard FB & DISC21 $\mu$AP \\
\midrule
PDQ & 0.579 & 0.520 & 0.495 & 0.495 \tiny[0.48, 0.50] & 0.000 & 0.295 \\
PDQ (dihedral) & 0.623 & 0.569 & 0.549 & 0.546 \tiny[0.54, 0.55] & 0.000 & 0.367 \\
aHash & 0.503 & 0.189 & 0.000 & 0.000 \tiny[0.00, 0.00] & 0.000 & 0.027 \\
dHash & 0.560 & 0.356 & 0.313 & 0.000 \tiny[0.00, 0.00] & 0.000 & 0.195 \\
pHash & 0.563 & 0.469 & 0.373 & 0.373 \tiny[0.00, 0.38] & 0.000 & 0.240 \\
pHash-256 & 0.589 & 0.520 & 0.494 & 0.492 \tiny[0.48, 0.50] & 0.001 & 0.302 \\
wHash & 0.425 & 0.134 & 0.000 & 0.000 \tiny[0.00, 0.00] & 0.000 & 0.024 \\
BlockMean & 0.562 & 0.338 & 0.337 & 0.000 \tiny[0.00, 0.00] & 0.001 & 0.069 \\
Marr–Hildreth & 0.586 & 0.455 & 0.362 & 0.314 \tiny[0.28, 0.35] & 0.001 & 0.292 \\
ColorMoment & 0.281 & 0.101 & 0.000 & 0.000 \tiny[0.00, 0.00] & 0.000 & 0.000 \\
Blockhash & 0.560 & 0.250 & 0.000 & 0.000 \tiny[0.00, 0.00] & 0.000 & 0.182 \\
ISCC Image-64 & 0.558 & 0.491 & 0.427 & 0.427 \tiny[0.31, 0.43] & 0.001 & 0.269 \\
ISCC Image-256 & 0.561 & 0.507 & 0.466 & 0.458 \tiny[0.45, 0.47] & 0.001 & 0.292 \\
DINOv2-S & 0.877 & 0.703 & 0.625 & 0.612 \tiny[0.60, 0.63] & 0.000 & 0.444 \\
DINOv2-B & 0.872 & 0.703 & 0.653 & 0.650 \tiny[0.64, 0.66] & 0.003 & 0.460 \\
DINOv2-S LSH-256 & 0.832 & 0.664 & 0.538 & 0.519 \tiny[0.50, 0.55] & 0.000 & 0.399 \\
OpenCLIP B/32 & 0.864 & 0.657 & 0.401 & 0.356 \tiny[0.33, 0.40] & 0.002 & 0.234 \\
SSCD R50\reftier & 0.929 & 0.763 & 0.725 & 0.712 \tiny[0.71, 0.73] & 0.007 & 0.795 \\
SSCD RX101\reftier & 0.942 & 0.787 & 0.763 & 0.703 \tiny[0.69, 0.77] & 0.007 & 0.823 \\
ISC21 1st\reftier & 0.919 & 0.815 & 0.762 & 0.755 \tiny[0.75, 0.77] & 0.004 & 0.917 \\
DINOHash-96\reftier & 0.769 & 0.533 & 0.237 & 0.237 \tiny[0.23, 0.24] & 0.000 & 0.212 \\
\bottomrule
\end{tabular}

\end{table*}


\paragraph{Where each class fails.} The heatmap explains the ranking. Every method, hash or learned, survives compression, rescaling and photometric changes almost perfectly. The hashes fail on geometry: a centred crop keeping 70\% of each side drops PDQ to \val{fp:image/PDQ/att/crop70/pct}\% Recall@1 and the ISCC Image-Code to \val{fp:image/ISCC-Image-64/att/crop70/pct}\%, while DINOv2-S keeps \val{fp:image/DINOv2-S/att/crop70/pct}\% and SSCD \val{fp:image/SSCD-mixup/att/crop70/pct}\%. A 15-degree rotation removes the hashes almost entirely (PDQ \val{fp:image/PDQ/att/rot15/pct}\%). PDQ's dihedral matching recovers flips and 90-degree rotations at no cost in precision (Table~\ref{tab:fp-paired}), which is why it leads the hashes. Two robustness properties are specific to ISCC: its uniform-border trim makes it the only hash that survives 25\% padding (\val{fp:image/ISCC-Image-64/att/pad25/pct}\%), while its $32\times32$ greyscale DCT is, apart from ColorMoment, the hash most affected by additive noise (\val{fp:image/ISCC-Image-64/att/noise0.05/pct}\% against \val{fp:image/PDQ/att/noise0.05/pct}\% for PDQ). Overlays are the hardest family for everyone: when the image is placed onto a screenshot template or onto another photograph at half size, the best method retrieves \val{fp:image/ISC21-1st/att/background0.5/pct}\% of the background overlays (ISC21) and \val{fp:image/ISC21-1st/att/screenshot/pct}\% of the screenshots, the best product-eligible one \val{fp:image/DINOv2-S/att/background0.5/pct}\% and \val{fp:image/DINOv2-S/att/screenshot/pct}\%, and every hash none. Such queries contain the registered image only as a region, and a single global descriptor cannot match it. Fig.~\ref{fig:fp-image-families}a shows the crop curve and Fig.~\ref{fig:fp-image-families}b the per-family means.

\begin{figure}[!t]
\centering
\includegraphics[width=\columnwidth]{fp_fig4_image_families}
\caption{(a) Recall@1 against the share of each side kept by a centred crop. (b) Mean Recall@1 over each transformation family for eight representative methods.}
\label{fig:fp-image-families}
\end{figure}

\paragraph{Size, speed and generalisation.} Fig.~\ref{fig:fp-image-tradeoffs} relates detection at the headline operating point to storage and extraction time. The keyed 256-bit projection of DINOv2-S keeps \val{fp:image/DINOv2-S-LSH256/TPR@pair@1e-07/pct}\% detection in 256 bits, against \val{fp:image/DINOv2-S/TPR@pair@1e-07/pct}\% for the 12\,288-bit float descriptor it is derived from. At PDQ's size it detects slightly fewer queries than dihedral PDQ at the strict threshold (\val{fp:image/PDQ-dihedral/TPR@pair@1e-07/pct}\%) but ranks the original first far more often (Recall@1 \val{fp:image/DINOv2-S-LSH256/R@1/pct}\% against \val{fp:image/PDQ-dihedral/R@1/pct}\%), because it survives crops and rotations. Panel~c shows that the ranking on ABO agrees with the ranking on our DISC21 subset, but that DISC21 separates the DISC21-trained descriptors far more sharply: SSCD and ISC21 reach $\mu$AP \val{fp:image/SSCD-mixup/disc} and \val{fp:image/ISC21-1st/disc} there, against \val{fp:image/DINOv2-S/disc} for DINOv2-S. DISC21's queries use the same augmentation library these models were trained with, and its references are photographs rather than product shots, so literature numbers on DISC21 overstate the gap on other content.

\begin{figure*}[!t]
\centering
\includegraphics[width=\textwidth]{fp_fig3_image_tradeoffs}
\caption{(a) Detection at pair-level FPR $10^{-7}$ against stored size per asset. (b) Against extraction time per 1024-pixel image (learned methods on one A10G GPU in batches of 64, hashes on one CPU thread; measured on an idle host). (c) $\mu$AP on ABO against $\mu$AP on the DISC21 subset.}
\label{fig:fp-image-tradeoffs}
\end{figure*}

\begin{table}[!t]
\centering
\footnotesize
\caption{Planned paired comparisons on the same image queries at pair-level FPR $10^{-7}$. The unit is the registered source with all its transformed queries: $p_{\mathrm{Holm}}$ from a sign-flip test on per-source differences in detections (Holm-adjusted over the eight comparisons); the interval of the TPR difference is a bootstrap over sources that also redraws the calibration negatives.}
\label{tab:fp-paired}
\setlength{\tabcolsep}{3pt}
\begin{tabular}{llcccrr}
\toprule
A & B & TPR$_A$ & TPR$_B$ & A$-$B [95\% CI] & sources & $p_{\mathrm{Holm}}$ \\
\midrule
SSCD R50 & DINOv2-S & 0.712 & 0.612 & 0.100 \tiny[0.084, 0.115] & 2{,}000 & $<10^{-3}$ \\
DINOv2-S & DINOv2-S LSH-256 & 0.612 & 0.519 & 0.093 \tiny[0.066, 0.106] & 2{,}000 & $<10^{-3}$ \\
PDQ & ISCC Image-64 & 0.495 & 0.427 & 0.068 \tiny[0.058, 0.175] & 2{,}000 & $<10^{-3}$ \\
PDQ & PDQ (dihedral) & 0.495 & 0.546 & -0.051 \tiny[-0.055, -0.050] & 2{,}000 & $<10^{-3}$ \\
ISCC Image-64 & ISCC Image-256 & 0.427 & 0.458 & -0.032 \tiny[-0.136, -0.020] & 2{,}000 & $<10^{-3}$ \\
SSCD R50 & ISC21 1st & 0.712 & 0.755 & -0.043 \tiny[-0.056, -0.029] & 2{,}000 & $<10^{-3}$ \\
DINOv2-S & OpenCLIP B/32 & 0.612 & 0.356 & 0.256 \tiny[0.220, 0.286] & 2{,}000 & $<10^{-3}$ \\
pHash & PDQ & 0.373 & 0.495 & -0.121 \tiny[-0.497, -0.112] & 2{,}000 & $<10^{-3}$ \\
\bottomrule
\end{tabular}

\end{table}


\subsection{False Matches, Variants and Registry Scale}
\label{sec:fp-res-fpr}
Thresholds fixed on the calibration negatives held out of sample: at the query-level 1\% threshold the held-out never-registered queries matched at rates between \val{fp:image/fprq/min}\% and \val{fp:image/fprq/max}\% across methods (Table~\ref{tab:fp-ci}). Fig.~\ref{fig:fp-scores}a shows why the learned descriptors can be operated at a strict threshold: their score to the true asset is well separated from the best score of a never-registered query, whereas for the hashes the two distributions meet near the threshold. Short codes also run out of resolution: at pair-level FPR $10^{-8}$ the 64-bit pHash and ISCC codes would need a Hamming distance below zero, so no threshold exists and their detection is \val{fp:image/pHash-64/TPR@pair@1e-08} and \val{fp:image/ISCC-Image-64/TPR@pair@1e-08}.

\begin{figure*}[!t]
\centering
\includegraphics[width=\textwidth]{fp_fig5_scores}
\caption{(a) Score distributions: best score of a never-registered query (grey) and score of an attacked query to its true asset (colour), pooled over all transformations; dashed line: threshold at pair-level FPR $10^{-7}$. (b) Rate at which another photograph of a registered product binds to that product at the 1\% threshold.}
\label{fig:fp-scores}
\end{figure*}

\paragraph{Different photographs of the same product.} A provenance registry must not bind a second photograph of a product to the manifest of the first. On \val{fp:dedup/n_hard} hard negatives the hashes almost never did, and the learned descriptors rarely did at the 1\% threshold: the highest rate was SSCD's \val{fp:image/SSCD-mixup/hard/pct}\% (Fig.~\ref{fig:fp-scores}b). Recoloured or re-mounted versions of the \emph{same} photograph are another matter. The \val{fp:dedup/n_hard_near_duplicate} hard negatives that our duplicate rule marked as variants bound to the registered image in most cases for the learned descriptors (e.g. 139 of 216 for SSCD-R50 and 96 for DINOv2-S) and rarely for the hashes (31 for PDQ, 3 for pHash).

\paragraph{Catalogue variants dominate the false-match budget.} ABO, like any real catalogue, reuses a design across products: the same artwork on phone cases of different models, the same ring with another stone. Table~\ref{tab:fp-sens} repeats the image evaluation with two other treatments of such variants. R@1 barely moves, but the calibrated operating points move a great deal: with no cleanup detection at the 1\% threshold falls to between \val{fp:sens/none/DINOv2-S-LSH256/TPR@query@0.01} and \val{fp:sens/none/SSCD-mixup/TPR@query@0.01} for every method (and to zero for pHash-256, which can then no longer set a threshold), because the threshold must then exclude near-copies of registered designs that are, in this benchmark, labelled as different assets. The choice of method matters less than whether the registry groups variants of one design under one asset.

\begin{table}[!t]
\centering
\footnotesize
\caption{Sensitivity of the image operating points to the treatment of catalogue duplicates. \emph{Any}: the rule fixed in advance (a pair counts as a duplicate if PDQ, SSCD or DINOv2-B says so); \emph{two}: two of the three must agree; \emph{none}: no cleanup, so recoloured or re-mounted versions of a registered design count as false matches. R@1 changes by at most a few points; the calibrated TPR does not.}
\label{tab:fp-sens}
\setlength{\tabcolsep}{3pt}
\begin{tabular}{lcccccc}
\toprule
 & \multicolumn{2}{c}{Any (primary)} & \multicolumn{2}{c}{Two of three} & \multicolumn{2}{c}{None} \\
Method & TPR$_{q,1\%}$ & TPR$_{p,10^{-7}}$ & TPR$_{q,1\%}$ & TPR$_{p,10^{-7}}$ & TPR$_{q,1\%}$ & TPR$_{p,10^{-7}}$ \\
\midrule
PDQ & 0.495 & 0.495 & 0.476 & 0.471 & 0.214 & 0.214 \\
PDQ (dihedral) & 0.549 & 0.546 & 0.523 & 0.523 & 0.223 & 0.223 \\
pHash-256 & 0.494 & 0.492 & 0.476 & 0.472 & 0.000 & 0.000 \\
ISCC Image-256 & 0.466 & 0.458 & 0.448 & 0.448 & 0.211 & 0.000 \\
DINOv2-S & 0.625 & 0.612 & 0.546 & 0.544 & 0.194 & 0.192 \\
DINOv2-S LSH-256 & 0.538 & 0.519 & 0.462 & 0.462 & 0.174 & 0.174 \\
DINOv2-B & 0.653 & 0.650 & 0.542 & 0.537 & 0.201 & 0.199 \\
OpenCLIP B/32 & 0.401 & 0.356 & 0.365 & 0.344 & 0.185 & 0.170 \\
SSCD R50\reftier & 0.725 & 0.712 & 0.601 & 0.512 & 0.236 & 0.217 \\
ISC21 1st\reftier & 0.762 & 0.755 & 0.704 & 0.665 & 0.213 & 0.212 \\
\bottomrule
\end{tabular}

\end{table}


\paragraph{Registry scale.} For $N$ registered assets and average pair-level false-match probability $p$, the expected number of false matches per negative query is $Np$ (Fig.~\ref{fig:fp-det}b). This count differs from the probability that a query has at least one false match, which is bounded above by $\min(1,Np)$. Detection falls as the target tightens (Fig.~\ref{fig:fp-det}a): from $10^{-7}$ to $10^{-8}$ SSCD-R50 loses \val{fp:image/SSCD-mixup/TPR@pair@1e-07/pct}\,$\to$\,\val{fp:image/SSCD-mixup/TPR@pair@1e-08/pct}\% and DINOv2-S \val{fp:image/DINOv2-S/TPR@pair@1e-07/pct}\,$\to$\,\val{fp:image/DINOv2-S/TPR@pair@1e-08/pct}\%. For a registry of ten million assets, an actual average per-reference false-match probability at most $10^{-9}$ for the relevant query population would bound the query-level false-match probability by 1\%. That operating point is below the empirically resolved range of this benchmark.

\paragraph{Does geometric verification close the gap?} A natural remedy is a second stage that checks the retrieved candidate against the stored original. We tested it on every positive and negative image query: the top-1 candidate of DINOv2-S or PDQ is accepted only if SIFT with RANSAC finds at least $k$ geometric inliers between the query and the registered image, with both the retrieval threshold and $k$ fixed on the calibration negatives (Table~\ref{tab:fp-verify}). For the evaluated DINOv2-S and PDQ pipelines, no detection improvement from verification was observed under a matched calibration false-binding constraint. With $k$ chosen so that the calibration false bindings do not exceed those of DINOv2-S alone at pair-level $10^{-7}$, the two-stage rule detected \val{fp:verify/DINOv2-S/query@0.01/k_match/tpr} of attacked queries against \val{fp:verify/DINOv2-S/one/tpr} for retrieval alone, and starting from the looser pair-level $10^{-4}$ threshold it detected only \val{fp:verify/DINOv2-S/pair@0.0001/k_match/tpr}. With the localization track's fixed $k=12$ it removed about four in ten held-out false bindings (\val{fp:verify/DINOv2-S/one/fp}$\,\to\,$\val{fp:verify/DINOv2-S/query@0.01/k12/fp} of \val{fp:verify/n_test}) at a cost of under a point of detection (\val{fp:verify/DINOv2-S/query@0.01/k12/tpr}), and PDQ behaved alike (\val{fp:verify/PDQ/one/fp}$\,\to\,$\val{fp:verify/PDQ/query@0.01/k12/fp}). Eliminating the remaining false bindings on the calibration set required hundreds of inliers (\val{fp:verify/DINOv2-S/query@0.01/k_zero/k} for DINOv2-S) and left \val{fp:verify/DINOv2-S/query@0.01/k_zero/tpr} detection. The false bindings that survive retrieval are, in large part, near-copies of registered designs, the catalogue variants of the sensitivity analysis, and they agree with the registered image geometrically as well as globally: a local-feature check confirms that two images share content, not that they are the same asset.

\begin{table}[!t]
\centering
\footnotesize
\caption{Second-stage verification of image matches. The top-1 candidate is accepted only if its retrieval score passes a threshold and SIFT+RANSAC finds at least $k$ geometric inliers between the query and the stored original. Both thresholds are fixed on the calibration negatives: matched, the smallest $k$ whose calibration false bindings do not exceed those of retrieval alone at pair-level $10^{-7}$; zero cal., the smallest $k$ no calibration negative reaches; and the localization track's $k=12$. FB: false bindings among the 26{,}820 held-out never-registered queries. TPR interval: bootstrap over sources.}
\label{tab:fp-verify}
\setlength{\tabcolsep}{3pt}
\begin{tabular}{llrccc}
\toprule
Method & Retrieval threshold & $k$ & TPR & FB & FB rate \\
\midrule
DINOv2-S & pair $10^{-7}$ & -- & 0.612 & 180 & 0.0067 \\
 & query 1\% & 9 (matched) & 0.609 \tiny[0.60, 0.61] & 124 & 0.0046 \\
 & query 1\% & 12 ($k=12$) & 0.604 \tiny[0.60, 0.61] & 104 & 0.0039 \\
 & query 1\% & 278 (zero cal.) & 0.206 \tiny[0.20, 0.22] & 0 & 0.0000 \\
 & pair $10^{-4}$ & 266 (matched) & 0.268 \tiny[0.25, 0.28] & 146 & 0.0054 \\
 & pair $10^{-4}$ & 12 ($k=12$) & 0.793 \tiny[0.79, 0.80] & 4{,}390 & 0.1637 \\
 & pair $10^{-4}$ & 1119 (zero cal.) & 0.023 \tiny[0.02, 0.03] & 0 & 0.0000 \\
\addlinespace
PDQ & pair $10^{-7}$ & -- & 0.495 & 267 & 0.0100 \\
 & query 1\% & 0 (matched) & 0.495 \tiny[0.49, 0.50] & 267 & 0.0100 \\
 & query 1\% & 12 ($k=12$) & 0.481 \tiny[0.48, 0.48] & 134 & 0.0050 \\
 & query 1\% & 644 (zero cal.) & 0.065 \tiny[0.06, 0.07] & 0 & 0.0000 \\
 & pair $10^{-4}$ & 167 (matched) & 0.299 \tiny[0.29, 0.31] & 212 & 0.0079 \\
 & pair $10^{-4}$ & 12 ($k=12$) & 0.559 \tiny[0.56, 0.56] & 2{,}565 & 0.0956 \\
 & pair $10^{-4}$ & 684 (zero cal.) & 0.066 \tiny[0.06, 0.07] & 0 & 0.0000 \\
\bottomrule
\end{tabular}

\end{table}


\begin{figure*}[!t]
\centering
\includegraphics[width=\textwidth]{fp_fig12_operating_points}
\caption{(a) Detection of attacked image queries against the pair-level false-positive rate at which the threshold is set. (b) Expected number of false matches per query as a function of registry size for fixed pair-level rates.}
\label{fig:fp-det}
\end{figure*}

\subsection{Audio}
\label{sec:fp-res-audio}
Fig.~\ref{fig:fp-audio} and Table~\ref{tab:fp-audio} report the audio track. As in the image track, the corpus contained duplicates: FMA carries some recordings under two track identifiers, which all three recording-level systems matched perfectly. We removed the \val{fp:audio/dedup/neg} negative sources that were copies of indexed tracks and counted a duplicate of a registered track as the same asset (\val{fp:audio/dedup/items} registered tracks had one), using the same any-of-three rule as for images. Without this step every method's calibrated detection falls by half or more (e.g. Chromaprint \val{fp:audio/Chromaprint/nodedup/tpr} instead of \val{fp:audio/Chromaprint/TPR@query@0.01}); we report the corrected values. With \val{fp:count/audio/neg_cal} calibration queries the pair-level targets below $10^{-6}$ are not resolved, so we compare at the query-level 1\% threshold.

The neural music fingerprint NMFP detects the most transformed queries (\val{fp:audio/NMFP-triplet/TPR@query@0.01/pct}\%), followed by Chromaprint (\val{fp:audio/Chromaprint/TPR@query@0.01/pct}\%) and audfprint (\val{fp:audio/audfprint/TPR@query@0.01/pct}\%); the ISCC Audio-Code reaches \val{fp:audio/ISCC-Audio-256/TPR@query@0.01/pct}\% at 256 bits and \val{fp:audio/ISCC-Audio-64/TPR@query@0.01/pct}\% at its default 64 bits, and the semantic CLAP embedding \val{fp:audio/CLAP/TPR@query@0.01/pct}\%. The families fail on different transformations (Fig.~\ref{fig:fp-audio}a). Codecs, resampling and moderate noise are harmless to all of them. Excerpts separate the systems that align in time from those that summarise a whole file: Chromaprint, audfprint and NMFP find a 5\,s excerpt of a 30\,s track in \val{fp:audio/Chromaprint/att/excerpt5/pct}\%, \val{fp:audio/audfprint/att/excerpt5/pct}\% and \val{fp:audio/NMFP-triplet/att/excerpt5/pct}\% of cases, while the ISCC Audio-Code, a single SimHash over the whole Chromaprint vector, finds \val{fp:audio/ISCC-Audio-64/att/excerpt5/pct}\%. Loud babble from other music (0\,dB) defeats the whole-file codes and CLAP, and most Chromaprint matches (\val{fp:audio/Chromaprint/att/babble_0db/pct}\%), while the landmark system audfprint (\val{fp:audio/audfprint/att/babble_0db/pct}\%), whose sparse spectral peaks survive mixing, and NMFP (\val{fp:audio/NMFP-triplet/att/babble_0db/pct}\%) keep most of them. Babble results depend on which background tracks are mixed in: the two runs of this benchmark drew different backgrounds (Section~\ref{sec:fp-method}), and NMFP's rate at 0\,dB was about half as high with the first set. Pitch and speed changes of a few percent defeat every recording-level fingerprint (Chromaprint \val{fp:audio/Chromaprint/att/pitch+2/pct}\% at $+2$ semitones, NMFP \val{fp:audio/NMFP-triplet/att/pitch+2/pct}\%), and only CLAP, which describes content rather than the recording, survives them (\val{fp:audio/CLAP/att/pitch+2/pct}\%) at the cost of failing under noise (\val{fp:audio/CLAP/att/noise_10db/pct}\% at 10\,dB SNR). The simulated re-recording is handled well by NMFP (\val{fp:audio/NMFP-triplet/att/rerecord/pct}\%) and audfprint (\val{fp:audio/audfprint/att/rerecord/pct}\%) and poorly by the others.

Speech behaves like music for the recording-level systems (Recall@1 on 10\,s excerpts: Chromaprint \val{fp:audio/Chromaprint/speech/att/excerpt10/pct}\% on speech and \val{fp:audio/Chromaprint/music/att/excerpt10/pct}\% on music), and ISCC does better on speech than on music excerpts (\val{fp:audio/ISCC-Audio-256/speech/att/excerpt10/pct}\% against \val{fp:audio/ISCC-Audio-256/music/att/excerpt10/pct}\% at 256 bits). NMFP is the only reference-tier audio method; it is GPL-3.0 and was trained on FMA, whose training tracks we could not exclude from our FMA-large corpus by identifier, so its lead may be partly in-distribution.

\begin{figure*}[!t]
\centering
\includegraphics[width=\textwidth]{fp_fig6_audio}
\caption{(a) Recall@1 of each audio method under each transformation (\val{fp:corpus/audio/reg} registered items, \val{fp:count/audio/nref} in the index). (b) Detection against the pair-level false-positive rate; unresolved targets are omitted.}
\label{fig:fp-audio}
\end{figure*}

\begin{table*}[!t]
\centering
\scriptsize
\caption{Audio retrieval over all attacks, and Recall@1 under selected attacks. TPR$_{q,1\%}$: detection at the threshold set for a 1\% query-level false-match rate on the calibration negatives; FPR$_q$ (test): the rate that threshold actually gave on the held-out negatives. TPR$_{p,10^{-6}}$: the strictest pair-level target that 4053 calibration queries against 6350 references resolve; stricter targets are not reported. Intervals: 95\% bootstrap over sources.}
\label{tab:fp-audio}
\setlength{\tabcolsep}{3pt}
\begin{tabular}{lccccccccc}
\toprule
Method & R@1 & $\mu$AP & TPR$_{q,1\%}$ & FPR$_{q}$ (test) & TPR$_{p,10^{-6}}$ & Music 10\,s exc. & Speech 10\,s exc. & Pitch +2 & Re-record \\
\midrule
Chromaprint & 0.805 & 0.817 & 0.704 \tiny[0.68, 0.72] & 0.016 \tiny[0.00, 0.04] & 0.688 \tiny[0.66, 0.71] & 1.000 & 1.000 & 0.006 & 0.451 \\
audfprint & 0.851 & 0.851 & 0.660 \tiny[0.53, 0.74] & 0.008 \tiny[0.00, 0.03] & 0.599 \tiny[0.46, 0.72] & 0.982 & 1.000 & 0.071 & 0.852 \\
ISCC Audio-64 & 0.662 & 0.581 & 0.409 \tiny[0.21, 0.41] & 0.017 \tiny[0.00, 0.03] & 0.212 \tiny[0.00, 0.41] & 0.247 & 0.657 & 0.004 & 0.183 \\
ISCC Audio-256 & 0.749 & 0.658 & 0.454 \tiny[0.31, 0.51] & 0.019 \tiny[0.00, 0.04] & 0.321 \tiny[0.24, 0.47] & 0.555 & 0.882 & 0.006 & 0.342 \\
CLAP & 0.707 & 0.446 & 0.386 \tiny[0.33, 0.45] & 0.017 \tiny[0.01, 0.04] & 0.328 \tiny[0.27, 0.41] & 0.957 & 0.725 & 0.910 & 0.228 \\
NMFP\reftier & 0.905 & 0.937 & 0.862 \tiny[0.82, 0.88] & 0.018 \tiny[0.00, 0.04] & 0.846 \tiny[0.79, 0.88] & 1.000 & 1.000 & 0.057 & 1.000 \\
\bottomrule
\end{tabular}

\end{table*}



\subsection{Video}
\label{sec:fp-res-video}
Fig.~\ref{fig:fp-video} and Table~\ref{tab:fp-video} report the video track. The \val{fp:count/video/neg_cal} calibration queries come from \val{fp:video/cal_sources} negative clips against \val{fp:count/video/nref} references, which resolves pair-level targets down to $10^{-4}$ but not below; Table~\ref{tab:fp-video} therefore reports the query-level 1\% threshold and the pair-level $10^{-4}$ target, and no stricter one. The query-level threshold did not deliver 1\% out of sample: on the held-out negatives it let through between \val{fp:video/fprq/min}\% and \val{fp:video/fprq/max}\% of never-registered queries, depending on the method. The threshold sits just above the sixth-highest of 500 calibration scores, so that five pass, and those queries are \val{fp:video/q_per_source} versions (the original and its transformations) of each of \val{fp:video/cal_sources} clips from a few films, whose clips resemble each other. The effective sample is the number of source clips, so the 1\% point is estimated imprecisely: the source-level 95\% intervals of the held-out rate reach up to \val{fp:video/fprq/cihi}\%. We report every video detection rate next to the false-match rate its threshold actually produced, and treat the video operating points as indicative. Frame-level learned descriptors are the most robust: SSCD and DINOv2-S reach Recall@1 of \val{fp:video/SSCD-mixup-mean/R@1/pct}\% and \val{fp:video/DINOv2-S-mean/R@1/pct}\% when their per-frame descriptors are averaged over the clip, and they are the only methods that survive 50\% crops, rotations and flips. Averaging, however, is also why their calibrated detection is lower than their Recall@1 suggests (\val{fp:video/SSCD-mixup-mean/TPR@query@0.01/pct}\% and \val{fp:video/DINOv2-S-mean/TPR@query@0.01/pct}\%): clips from the same film average to similar vectors, which raises the threshold. The ISCC Video-Code and TMK+PDQF are less robust but detect more at the 1\% threshold (\val{fp:video/ISCC-Video-256/TPR@query@0.01/pct}\% and \val{fp:video/TMK+PDQF/TPR@query@0.01/pct}\%), at the price of more held-out false matches (\val{fp:video/ISCC-Video-256/FPRq@query@0.01/pct}\% and \val{fp:video/TMK+PDQF/FPRq@query@0.01/pct}\%, against \val{fp:video/SSCD-mixup-mean/FPRq@query@0.01/pct}\% and \val{fp:video/DINOv2-S-mean/FPRq@query@0.01/pct}\% for the averaged learned descriptors); with 20 calibration clips these differences in false-match rate are within the intervals of Table~\ref{tab:fp-video}. Picture-in-picture defeats every method, as the overlay attacks did for images.

\begin{figure*}[!t]
\centering
\includegraphics[width=\textwidth]{fp_fig7_video}
\caption{(a) Recall@1 of each video method under each transformation (\val{fp:corpus/video/reg} registered clips, \val{fp:count/video/nref} in the index). (b) Detection against the pair-level false-positive rate; unresolved targets are omitted.}
\label{fig:fp-video}
\end{figure*}

\begin{table*}[!t]
\centering
\scriptsize
\caption{Video retrieval over all attacks, and Recall@1 under selected attacks. TPR$_{q,1\%}$: detection at the threshold set for a 1\% query-level false-match rate on the calibration negatives; FPR$_q$ (test): the rate that threshold actually gave on the held-out negatives. TPR$_{p,10^{-4}}$: the strictest pair-level target that 500 calibration queries against 443 references resolve; stricter targets are not reported. Intervals: 95\% bootstrap over sources.}
\label{tab:fp-video}
\setlength{\tabcolsep}{3pt}
\begin{tabular}{lccccccccc}
\toprule
Method & R@1 & $\mu$AP & TPR$_{q,1\%}$ & FPR$_{q}$ (test) & TPR$_{p,10^{-4}}$ & Crop 50\% & PiP & Inserted & Screen rec. \\
\midrule
vPDQ & 0.628 & 0.617 & 0.577 \tiny[0.56, 0.59] & 0.032 \tiny[0.00, 0.08] & 0.609 \tiny[0.57, 0.62] & 0.000 & 0.000 & 0.342 & 0.000 \\
TMK+PDQF & 0.778 & 0.723 & 0.631 \tiny[0.61, 0.66] & 0.066 \tiny[0.00, 0.16] & 0.635 \tiny[0.63, 0.70] & 0.017 & 0.000 & 0.583 & 0.992 \\
ISCC Video-64 & 0.731 & 0.700 & 0.620 \tiny[0.61, 0.65] & 0.062 \tiny[0.00, 0.16] & 0.634 \tiny[0.62, 0.66] & 0.025 & 0.000 & 0.333 & 0.817 \\
ISCC Video-256 & 0.801 & 0.749 & 0.652 \tiny[0.63, 0.67] & 0.068 \tiny[0.00, 0.17] & 0.664 \tiny[0.66, 0.69] & 0.108 & 0.008 & 0.575 & 0.958 \\
DINOv2-S frames & 0.872 & 0.691 & 0.437 \tiny[0.41, 0.58] & 0.062 \tiny[0.01, 0.15] & 0.523 \tiny[0.45, 0.63] & 0.592 & 0.125 & 0.525 & 0.817 \\
DINOv2-S mean & 0.947 & 0.831 & 0.547 \tiny[0.51, 0.73] & 0.032 \tiny[0.00, 0.16] & 0.653 \tiny[0.56, 0.77] & 0.917 & 0.200 & 0.958 & 0.950 \\
SSCD frames\reftier & 0.921 & 0.806 & 0.502 \tiny[0.48, 0.56] & 0.032 \tiny[0.00, 0.10] & 0.535 \tiny[0.50, 0.69] & 0.917 & 0.067 & 0.658 & 0.983 \\
SSCD mean\reftier & 0.958 & 0.893 & 0.543 \tiny[0.52, 0.78] & 0.036 \tiny[0.00, 0.14] & 0.697 \tiny[0.55, 0.83] & 0.983 & 0.108 & 0.950 & 0.992 \\
\bottomrule
\end{tabular}

\end{table*}


\subsection{Partial Edits and Localization}
\label{sec:fp-res-edits}
Partial edits cost the global descriptors little in ranking but much in calibrated detection. SSCD-R50 still ranks the original first for \val{fp:editret/SSCD-mixup/R@1}\ of the \val{fp:edit/n_edits} edited images and DINOv2-S for \val{fp:editret/DINOv2-S/R@1}, while PDQ does for \val{fp:editret/PDQ/R@1} and the ISCC Image-Code for \val{fp:editret/ISCC-Image-64/R@1}. At the 1\% threshold, however, an edit covering a quarter of the image leaves SSCD's score above threshold in only \val{fp:editret/SSCD-mixup/query@0.01/area/0.25} of cases, against \val{fp:editret/SSCD-mixup/query@0.01/area/0.01} for 1\% edits (Fig.~\ref{fig:fp-loc}c). A fingerprint that tolerates an edit well enough to find the original is, by design, not a tamper detector.

Once the original is found, the edit can be localized against it (Table~\ref{tab:fp-loc}, Fig.~\ref{fig:fp-loc}). Table~\ref{tab:fp-loc} is an oracle measurement: each edited query is compared with its true registered original, whether or not a fingerprint would have retrieved it. SIFT alignment succeeded for \val{fp:loc/sift/aligned}\% of the evaluated edits and DISK+LightGlue for \val{fp:loc/lightglue/aligned}\%. The simplest map, a blurred pixel difference after alignment, gave the best F1 at the calibrated 1\% pixel false-positive rate (\val{fp:loc/sift/pixel/f1} with SIFT), and LPIPS the highest AUC (\val{fp:loc/sift/lpips/auc}); the fused map was the most robust to the 90\% crop, which LightGlue aligns better than SIFT (fused F1 \val{fp:loc/lightglue/fused/post/crop90/f1} against \val{fp:loc/sift/fused/post/crop90/f1}). Text replacement and recolouring were the easiest edits (pixel F1 \val{fp:loc/sift/pixel/type/textreplace/f1} and \val{fp:loc/sift/pixel/type/recolor/f1}) and learned inpainting the hardest (\val{fp:loc/sift/pixel/type/inpaint-lama/f1}), because a good inpainting of a textureless background changes few pixels. PDQ computed per tile could not localize at all at a 1\% false-positive rate (F1 \val{fp:loc/sift/pdq/f1}): JPEG re-encoding alone flips tile hashes. End to end, the original must first be retrieved. Counting an edit whose original a fingerprint does not rank first as a localization failure, the SIFT pixel-difference F1 over the \val{fp:loc/e2e/sift/n} evaluated edits falls from \val{fp:loc/sift/pixel/f1} (oracle) to \val{fp:loc/e2e/sift/SSCD-mixup/pixel/f1} with SSCD-R50 retrieval, \val{fp:loc/e2e/sift/DINOv2-S/pixel/f1} with DINOv2-S, \val{fp:loc/e2e/sift/PDQ/pixel/f1} with PDQ and \val{fp:loc/e2e/sift/ISCC-Image-64/pixel/f1} with the ISCC Image-Code, which ranked the original first for \val{fp:loc/e2e/sift/SSCD-mixup/retrieved}, \val{fp:loc/e2e/sift/DINOv2-S/retrieved}, \val{fp:loc/e2e/sift/PDQ/retrieved} and \val{fp:loc/e2e/sift/ISCC-Image-64/retrieved}\% of them. These figures use top-1 retrieval without a threshold; requiring the calibrated threshold as well lowers them further, by the gap between Recall@1 and TPR for edited queries reported above.

\begin{table}[!t]
\centering
\footnotesize
\caption{Localization of partial edits against the registered original (oracle retrieval; evaluation half of the sources). F1 at the threshold fixed for a 1\% pixel false-positive rate on unedited, post-processed originals; best F1 chooses the threshold per map with hindsight.}
\label{tab:fp-loc}
\setlength{\tabcolsep}{3pt}
\begin{tabular}{lcccccc}
\toprule
 & \multicolumn{3}{c}{SIFT + RANSAC} & \multicolumn{3}{c}{DISK + LightGlue} \\
Map & AUC & F1 & best F1 & AUC & F1 & best F1 \\
\midrule
Pixel difference & 0.944 & 0.75 & 0.83 & 0.942 & 0.72 & 0.82 \\
SSIM & 0.927 & 0.57 & 0.74 & 0.924 & 0.59 & 0.71 \\
LPIPS & 0.975 & 0.64 & 0.79 & 0.973 & 0.61 & 0.78 \\
DINOv2 patches & 0.958 & 0.53 & 0.70 & 0.952 & 0.52 & 0.69 \\
PDQ tiles & 0.866 & 0.04 & 0.50 & 0.866 & 0.02 & 0.49 \\
Fused & 0.967 & 0.65 & 0.81 & 0.965 & 0.69 & 0.79 \\
\bottomrule
\end{tabular}

\end{table}


\begin{figure*}[!t]
\centering
\includegraphics[width=\textwidth]{fp_fig8_localization}
\caption{(a) Examples of each edit type: registered original, edited query, ground-truth region, and three difference maps after SIFT alignment. (b) Pixel F1 at the calibrated threshold against the edited area, per map. (c) Share of edited queries whose original is retrieved above the query-level 1\% threshold, against edited area, for the eight representative methods of Fig.~\ref{fig:fp-image-families}.}
\label{fig:fp-loc}
\end{figure*}

For time-based media, localization means finding where a query sits in the registered item. Chromaprint recovered the offset of 10\,s and 5\,s excerpts and of re-recordings within 0.5\,s in \val{fp:tloc/audio/chromaprint/_all/within05}\% of cases (median error \val{fp:tloc/audio/chromaprint/_all/median_ms}\,ms) and audfprint in \val{fp:tloc/audio/audfprint/_all/within05}\% (median \val{fp:tloc/audio/audfprint/_all/median_ms}\,ms); its failures are mostly queries it did not match at all. For video, frame-level SSCD and DINOv2-S placed every 2.5\,s excerpt within 1\,s of its true position, and recovered the inserted interval of the insertion attack with a mean interval IoU of \val{fp:tloc/video/SSCD-mixup/insert/iou} and \val{fp:tloc/video/DINOv2-S/insert/iou}.

\subsection{Security}
\label{sec:fp-res-security}
\paragraph{Evasion.} Every learned descriptor could be evaded by a white-box attacker with an imperceptible perturbation (Fig.~\ref{fig:fp-security}a): at a budget of 2 grey levels (PSNR \val{fp:sec/whitebox/DINOv2-S/2/psnr}\,dB) projected gradient descent pushed \val{fp:sec/whitebox/DINOv2-S/2/pct7}\% of the registered images below the $10^{-7}$ threshold for DINOv2-S, and equally for SSCD, ISC21, OpenCLIP and the keyed projection, whose soft pre-sign values carry the gradient. The DCT hashes resisted longer but not for long: through differentiable surrogates the attack evaded PDQ for \val{fp:sec/surrogate/PDQ/2/pct7}\%, \val{fp:sec/surrogate/PDQ/4/pct7}\% and \val{fp:sec/surrogate/PDQ/8/pct7}\% of images at 2, 4 and 8 grey levels, pHash for \val{fp:sec/surrogate/pHash-64/2/pct7}\% already at 2, and the ISCC Image-Code for \val{fp:sec/surrogate/ISCC-Image-64/2/pct7}\% at 2. A black-box attacker who crafts the perturbation on an ensemble of DINOv2-S, SSCD and OpenCLIP (Fig.~\ref{fig:fp-security}b, budget 8) evades every DINOv2 variant and SSCD-R50 but transfers poorly to ISC21 (\val{fp:sec/transfer/8/ISC21-1st/pct}\%) and not at all to PDQ and pHash-256; the ISCC codes are evaded for about half of the images at every budget (\val{fp:sec/transfer/2/ISCC-Image-64/pct}\% at 2 grey levels), again through the border-trimming step, which small perturbations of a white background switch on or off.

\paragraph{Forgery.} Binding an unrelated image to a registered asset is harder but, for the learned copy detectors, practical (Fig.~\ref{fig:fp-security}c): at 4 grey levels the attack made \val{fp:sec/collision/ISC21-1st/4/pct}\% of source images the top match of their target in the full registry above threshold for ISC21 and \val{fp:sec/collision/SSCD-mixup/4/pct}\% for SSCD-R50, rising to \val{fp:sec/collision/ISC21-1st/8/pct}\% and \val{fp:sec/collision/SSCD-mixup/8/pct}\% at 8. DINOv2-S (\val{fp:sec/collision/DINOv2-S/8/pct}\% at 8) and OpenCLIP (\val{fp:sec/collision/OpenCLIP-B32/8/pct}\%) were harder to steer, and the hashes nearly impossible within these budgets (PDQ \val{fp:sec/collision/PDQ/16/pct}\% and ISCC \val{fp:sec/collision/ISCC-Image-64/16/pct}\% at 16 grey levels, where the PSNR is already \val{fp:sec/collision/PDQ/16/psnr}\,dB). The copy-detection descriptors, which generalise best across benign transformations, are thus also the easiest to forge: their similarity is smooth in the pixels, which is what both robustness and gradient steering exploit.

\paragraph{Inversion.} A soft-binding value is published inside the manifest, so we asked how much of the image it gives away. A decoder trained on the descriptors of the \val{fp:sec/inv/ntrain} distractors reconstructed a $64\times64$ thumbnail of each of \val{fp:sec/inv/ntest} registered images from its descriptor alone (Fig.~\ref{fig:fp-inversion}). Reconstructions from the float descriptors were recognisable in colour, layout and object class (LPIPS \val{fp:sec/inv/DINOv2-S/lpips} for DINOv2-S against \val{fp:sec/inv/mean-image/lpips} for the mean image), and using them as queries retrieved the true image among all \val{fp:sec/inv/ntest} in \val{fp:sec/inv/OpenCLIP-B32/reid/pct}\% (OpenCLIP), \val{fp:sec/inv/DINOv2-S/reid/pct}\% (DINOv2-S) and \val{fp:sec/inv/SSCD-mixup/reid/pct}\% (SSCD) of cases, 40 to 95 times chance. The binary hashes leaked less but still identified images above chance (PDQ \val{fp:sec/inv/PDQ/reid/pct}\%). The keyed projection behaved as intended: an attacker who knows the key still identifies images at about 40 times chance (\val{fp:sec/inv/DINOv2-S-LSH256/reid/pct}\%, half the rate from the float DINOv2-S descriptor), but a decoder trained with any other key reconstructs only the mean image (LPIPS \val{fp:sec/inv/DINOv2-S-LSH256-wrongkey/lpips}) and identifies images at chance (\val{fp:sec/inv/DINOv2-S-LSH256-wrongkey/reid/pct}\%).

\begin{figure*}[!t]
\centering
\includegraphics[width=\textwidth]{fp_fig10_inversion}
\caption{Inverting stored fingerprints. (a) Registered images (top) and thumbnails reconstructed from their descriptors by a decoder trained on distractor descriptors; bottom rows: the keyed 256-bit projection decoded by an attacker with the right and with a wrong key. (b) Share of reconstructions that retrieve their own original among the \val{fp:sec/inv/ntest} registered images (95\% Clopper--Pearson intervals); dashed: chance.}
\label{fig:fp-inversion}
\end{figure*}


\begin{figure*}[!t]
\centering
\includegraphics[width=\textwidth]{fp_fig9_security}
\caption{Adversarial robustness on \val{fp:sec/n} registered ABO images per setting. (a) Evasion by white-box PGD (learned descriptors) and surrogate-gradient PGD (DCT hashes) against the $L_\infty$ budget, judged by the real extractor at pair-level FPR $10^{-7}$. (b) Evasion by perturbations crafted on a DINOv2-S, SSCD-R50 and OpenCLIP ensemble at a budget of 8 grey levels, judged at the pair-level $10^{-6}$ threshold the attack ran against; n/a: the method has no reachable operating point. (c) Targeted forgery: share of unrelated images that become their target's top match above the $10^{-6}$ threshold.}
\label{fig:fp-security}
\end{figure*}

\paragraph{Audio and video.} The same pattern holds in the other modalities. Audio attacks were judged, like the image attacks, at the pair-level $10^{-6}$ threshold of each method, taken from the first audio retrieval run; those thresholds differ from the re-run reported in Section~\ref{sec:fp-res-audio} by at most \val{fp:sec/audio/thdelta}, and re-scoring every CLAP attack at the re-run thresholds changes \val{fp:sec/audio/rescore_changed} outcomes. Video attacks were judged at a different threshold, described below. For audio, projected gradient descent on the CLAP embedding of \val{fp:sec/audio/n} registered tracks evaded CLAP for every track at every budget we tried, down to a target signal-to-noise ratio of 40\,dB (achieved median \val{fp:sec/audio/snr/40}\,dB). The perturbation did not transfer to the recording-level fingerprints: Chromaprint matched every attacked track at all budgets, and the ISCC Audio-Code and audfprint lost \val{fp:sec/audio/pgd/20/ISCC-Audio-64/pct}\% and \val{fp:sec/audio/pgd/20/audfprint/pct}\% only at the loudest budget (20\,dB target), where the added signal acts as noise rather than as a targeted attack (\val{fp:sec/audio/pgd/40/ISCC-Audio-64/pct}\% and \val{fp:sec/audio/pgd/40/audfprint/pct}\% at 40\,dB). An attacker does not need gradients against those systems, however: a search over pitch and tempo changes found the smallest change that broke each match, and the median was \val{fp:sec/audio/mag/Chromaprint/pitch} semitones or a tempo factor of \val{fp:sec/audio/mag/Chromaprint/tempo} for Chromaprint and \val{fp:sec/audio/mag/audfprint/pitch} semitones or \val{fp:sec/audio/mag/audfprint/tempo} for audfprint and the ISCC Audio-Code, where 0.25 semitones and 1.02 are the smallest steps of the grid. Only CLAP needed audible changes (\val{fp:sec/audio/mag/CLAP/pitch} semitones, tempo \val{fp:sec/audio/mag/CLAP/tempo}). Every audio fingerprint in this study is therefore evaded by a small change, either by gradients (CLAP) or by a pitch shift (all others).

For video, frame-wise projected gradient descent against DINOv2-S and SSCD on \val{fp:sec/video/n} registered clips evaded both frame-level neural descriptors, in sequence and mean-pooled form, for every clip at 4 grey levels, and transferred to none of the hashes: vPDQ, TMK+PDQF and both ISCC Video-Codes matched every attacked clip at 4 and 8 grey levels (judged at each method's query-level 1\% threshold). The video calibration set does not resolve a pair-level $10^{-6}$ target, so every video attack used that threshold, which let through \val{fp:video/fprq/min}\% to \val{fp:video/fprq/max}\% of held-out negative queries (Section~\ref{sec:fp-res-video}). The video results therefore do not establish attack performance at a pair-level false-match rate. Each hash's pair-level $10^{-4}$ threshold is lower than its query-level threshold, so the hashes' matches hold at that target too; the per-clip scores of the neural descriptors were not kept, so their evasion was not re-scored at it. Simple edits again did what gradients could not. A centred crop keeping 90\% of the width and height, the first step of the grid, broke the median match for vPDQ, TMK+PDQF, both ISCC Video-Codes and the sequence-matched DINOv2-S; only SSCD and mean-pooled DINOv2-S held at 90\% and broke at a median of \val{fp:sec/video/magnitude/SSCD-mixup-seq/crop_keep_median}. Playback speed-up was the reverse: TMK+PDQF and the 256-bit ISCC Video-Code still matched all \val{fp:sec/video/magnitude/TMK+PDQF/speed_never} clips at twice the speed (the 64-bit code \val{fp:sec/video/magnitude/ISCC-Video-64/speed_never}), because dropping frames leaves the remaining frames' hashes intact, while the sequence-matched neural descriptors broke at a median factor of \val{fp:sec/video/magnitude/DINOv2-S-seq/speed_median}. No video fingerprint survived both a small crop and a white-box attack.




\subsection{Determinism and Cost}
\label{sec:fp-res-cost}
\paragraph{Determinism across hosts.} A registry compares a fingerprint computed at signing time with one computed at lookup time, possibly on different hardware. We fingerprinted the same \val{fp:stab/n} PNG files (100 registered originals and their JPEG-75 copies, byte-identical on both hosts) on the AMD EPYC and NVIDIA A10G host and on an Apple M5 Pro laptop (Fig.~\ref{fig:fp-stability}). PDQ, the imagehash DCT and difference hashes, BlockMean, Blockhash and the ISCC codes were bit-identical on every image, and the float descriptors agreed to $1-\cos \le$ \val{fp:stab/SSCD-large/cosmax}. Two classical hashes were not deterministic: the Marr--Hildreth hash differed on \val{fp:stab/MarrHildreth/differ/pct}\% of images by up to \val{fp:stab/MarrHildreth/bitsmax} of its 576 bits, enough to push \val{fp:stab/MarrHildreth/below} of the \val{fp:stab/n} self-comparisons below its own match threshold, and wHash differed on \val{fp:stab/wHash-64/differ/pct}\% by up to \val{fp:stab/wHash-64/bitsmax} of 64 bits. Both depend on floating-point filtering whose result differs between processor architectures; a service that fingerprints on one platform and verifies on another should avoid them or pin the implementation.

\paragraph{Extraction cost.} We timed extraction separately on the otherwise idle A10G host (AMD EPYC 7R32), because the timings recorded during the benchmark include contention from dozens of concurrent workers. Each image method fingerprinted the same 256 registered originals at 1024 pixels on the long side, CPU methods on a single thread and learned methods on the GPU in batches of 64 including the host-to-device copy. The hashes cost between \val{fp:lat/image/BlockMean} (BlockMean) and \val{fp:lat/image/wHash-64}\,ms (wHash) per image on one CPU thread, with PDQ at \val{fp:lat/image/PDQ}\,ms and the ISCC Image-Code at \val{fp:lat/image/ISCC-Image-64}\,ms; the learned descriptors cost \val{fp:lat/image/OpenCLIP-B32}--\val{fp:lat/image/ISC21-1st}\,ms on the GPU (DINOv2-S \val{fp:lat/image/DINOv2-S}, SSCD-R50 \val{fp:lat/image/SSCD-mixup}, ISC21 \val{fp:lat/image/ISC21-1st}). None of these costs matters for a signing service: on the same on-demand instance, fingerprinting a million images costs about US\$\val{fp:lat/image/PDQ/usd1m} with PDQ on its eight vCPUs, US\$\val{fp:lat/image/DINOv2-S/usd1m} with DINOv2-S and US\$\val{fp:lat/image/ISC21-1st/usd1m} with ISC21 on its GPU, before decoding. Storage differs more: PDQ, the ISCC codes and the keyed projection need 8--32 bytes per asset, a float DINOv2-S descriptor 1.5\,kB. Audio fingerprints ran at \val{fp:lat/audio/audfprint}--\val{fp:lat/audio/ISCC-Audio-256}\,ms per second of audio (CLAP \val{fp:lat/audio/CLAP}\,ms on the GPU), and video fingerprints at \val{fp:lat/video/TMK+PDQF}\,ms (TMK+PDQF) to \val{fp:lat/video/vPDQ}\,ms (vPDQ) per second of video on one CPU thread and \val{fp:lat/video/DINOv2-S-seq}\,ms for frame-level DINOv2-S, dominated in every case by decoding and frame sampling. We did not time registry search at scale; exact search over our \val{fp:count/image/nref}-item image index was not the bottleneck at any point.


\begin{figure}[!t]
\centering
\includegraphics[width=\columnwidth]{fp_fig13_stability}
\caption{Share of binary image codes that differ between an x86/A10G host and an arm64/M5 Pro laptop on byte-identical inputs. Float descriptors differ only in rounding ($1-\cos \le 2\times10^{-4}$).}
\label{fig:fp-stability}
\end{figure}

